


Le schéma de base de données suivant sera utilisé pour l'ensemble du
sujet. Il permet de gérer une base sur le guide du Routard.

\begin{tabular}{ll}
{\sc Guide} &(\textbf{id}, titre, année)\\
{\sc Ville} &(\textbf{id}, nomVille, \textit{idGuide}, pays)\\
{\sc Site} &(\textbf{id}, nomSite, type, \textit{idVille}, GPS)\\
{\sc Routard} &(\textbf{id}, nom, prénom, pays)\\
{\sc Commentaire} &(\textbf{\textit{idSite}, \textit{idRoutard}, date}, note, commentaire)
\end{tabular}\\

Tous les attributs \textbf{en gras}
sont des clés primaires, ceux \textit{en italiques} sont des clés étrangères. Notez
qu'une clé étrangère peut faire partie d'un clé primaire.

Le {\sc guide} est l'édition d'un ``Guide du Routard'' 
à une année donnée. Un guide est consacré à un ensemble de villes contenant des sites.
Le type de {\sc Site} peut être : Hôtel, Restaurant ou PdI (``Point d'Intérêt'').
La table {\sc Commentaire} permet d'associer le commentaire d'un routard avec 
un site à une date donnée. 

\section{Conception (5 points)}

\begin{enumerate}
    \item Si je connais un site, est-ce que je sais dans quel guide il est décrit (penser dépendance fonctionnelle
    et transitivité)\,? Expliquez.
    \item Un routard peut-il commenter un même site plusieurs fois\,? Si oui à quelle condition\,?
     \item Donnez le schéma entité-association correspondant aux relations \textsc{Guide}, \textsc{Ville},
          \textsc{Site}.
     \item Donnez la commande SQL de création de la table \textsc{Commentaire}.
     \item Supposons que l'on  modifie le schéma pour identifier un site par la paire (\textbf{idVille, nomSite})
     en supprimant l'attribut \textbf{id}..
             Est-ce que cela correspond d'après vous à une modélisation par entité faible ou
             par réification\,? Que faut-il modifier d'autre au schéma initial\,?
\end{enumerate}

\begin{solution}
\begin{enumerate}
  \item Oui, par transitivité.
  \item Oui, à des dates différentes
\end{enumerate}

\end{solution}
\section{SQL (8 points)}

\begin{enumerate}

\item Noms des routards ayant  écrit au moins un commentaire sur le Parthénon (qui
est un nom de site).
\begin{solution}
\begin{minted}{sql}
select nom
from Routard as r, Commentaire as c, Site as s
where s.nomSite='Parthénon'
and c.idSite=s.id
and c.idRoutard = r.id
\end{minted}
\end{solution}


\item Indiquez ce que fait  la requête suivante et donnez une syntaxe équivalente, \emph{sans imbrication}.

\begin{minted}{sql}
select id, titre
from Guide as g
where id in (select idGuide from Ville as v where nomVille='Venise')
\end{minted}


\begin{solution}
\begin{minted}{sql}
select g.id, g.titre
from Guide as g, Ville as v
where g.id=v.idGuide
and nomVille='Venise'
\end{minted}
\end{solution}


\item Noms des routards ayant  écrit au moins un commentaire avec une note inférieure ou égale à 2 
   sur un site de leur propre pays.

\begin{solution}
\begin{minted}{sql}
select nom
from Routard as r, Commentaire as c, Site as s, Ville as v
where c.note < 3
and c.idSite=s.id
and v.id=s.idVille
and v.pays=r.pays
and c.idRoutard = r.id
\end{minted}
\end{solution}


\item Quels sont les noms des sites sur lesquels aucun routard français n'a écrit de commentaire? 

\begin{solution}
\begin{minted}{sql}
select nomSite
from Site as s
where not exists (select '' 
            from Routard as r, Commentaire as c
            where c.idRoutard = r.id
            and r.pays='France'
            and s.id=c.idSite)
\end{minted}
\end{solution}

\item Quels sont les sites (donner le nom) qui ont une note moyenne strictement supérieure à 4?

\begin{solution}
\begin{minted}{sql}
select nomSite
from Site as s, Commentaire as c
where s.id=c.idSite
group by s.id
having avg (note) >= 4
\end{minted}
\end{solution}


\item Quelles paires de sites ont le même nom, mais dans des villes différentes. Donner le nom du site, et
     les noms des deux villes, comme, par exemple, ('Panthéon, 'Rome', 'Paris')
 

\begin{solution}
\begin{minted}{sql}
select s1.nomSite as 'nomSite', v1.nom as 'Ville1', v2.nom as 'Ville2'
from Site as s1, Site as s2, Ville as v1, Ville as v2
where s1.idVille=v1.id
and s2.idVille=v2.id
and not (v1.id = v2.id)
and s1.nomSite = s2.nomSite
\end{minted}
\end{solution}

\item  Quelle définition de la notion de vue parmi les suivantes vous semble correcte ?

    \begin{enumerate}
        \item  C’est un moyen de pré-calculer un résultat.
       \item  C’est la création d'une seconde base dérivée de la première par exécution de requêtes SQL.
       \item C’est un moyen de prédéfinir des requêtes en les nommant.
\end{enumerate}


\begin{solution}
Réponse c. Les vues en SQL sont des tables intentionnelles, leur contenu n’est calculé qu’au moment où on les interroge. 
\end{solution}

\end{enumerate}

\section{Algèbre (3 points)}

\begin{enumerate}
   \item L’évaluation d’une  de ces opérations algébriques demande l’élimination des doublons, laquelle :
 \begin{enumerate}
        \item  La projection
       \item  Le renommage
       \item La sélection
\end{enumerate}


\item Donnez l'expression algébrique pour, au choix,  une des 3 premières requêtes de la section précédente

\item Comment interprétez-vous l'expression algébrique suivante\,? Quel résultat donne-t-elle\,?
$$\pi_{id}(Guide) - \pi_{idGuide} (\sigma_{pays = '\rm{Italie}'}(Ville))$$
\begin{solution}
Les identifiants des guides qui ne parlent d'aucune ville en Italie.
\end{solution}
\end{enumerate}

\section{Programmation et transactions (4 points)}

\begin{enumerate}

  \item Quel est l’intérêt d’associer un langage de programmation à SQL ?

\begin{enumerate}
  \item  SQL étant déclaratif, il faut bien à un moment introduire un langage pour programmer l’évaluation des requêtes

  \item SQL ne permet pas d’effectuer des boucles ou des tests, un langage de programmation le complète donc
  \item On dispose ainsi de deux moyens d’effectuer des requêtes, soit par SQL, soit avec un langage généraliste
\end{enumerate}


\begin{solution}
Réponse 2. L’évaluation des requêtes est du ressort du système, 
pas du développeur d’application. En ce qui concerne la troisième affirmation, il est infiniment plus facile d’interroger une base par SQL que d’écrire un programme pour chaque requête.
\end{solution}

\item Qu’est-ce qu’un curseur ?

\begin{enumerate}
  \item  Un tableau contenant le résultat d’une requête
  \item Un mécanisme permettant de parcourir le résultat d’une requête
  \item Une variable indiquant combien de nuplets on souhaite récupérer après exécution d’une requête
\end{enumerate}

\begin{solution}
Réponse 2. un curseur n’est pas une structure contenant le résultat. Ce résultat peut d’ailleurs ne jamais exister intégralement, 
mais être calculé, puis oublié au fur et à mesure.
\end{solution}

\item À quel moment doit-on effectuer un \texttt{commit}? (1 point)

\begin{enumerate}
  \item  Dès que la base est arrivée à un état cohérent
  \item  Après chaque mise à jour
  \item  À intervalles périodiques (par exemple toutes les 5 mns)
\end{enumerate}


\begin{solution}
Réponse 1
\end{solution}

\item Quel est l'intérêt du mode \texttt{serializable} par rapport aux autres niveaux d'isolation? Citer au moins un inconvénient. (1 point)

\begin{solution}
Garantie de l'isolation totale et de la correction transactionnelle; mais risque de deadlock.
\end{solution}



\end{enumerate}
